\chapter{Clase 18: Geometría de las cónicas (parte 1)}

Esta clase inicia el estudio geométrico de las cónicas desde la perspectiva de la mecánica celeste. Después de haber deducido la ecuación polar de la trayectoria relativa, el siguiente paso es comprender en qué sistema de referencia esa ecuación adopta su forma más simple y cómo se transforma al sistema de referencia del observador. El objetivo es establecer con claridad el \textbf{sistema natural de la cónica} y desarrollar la herramienta algebraica básica para relacionarlo con un sistema externo: la \textbf{rotación plana}.

\section{Introducción y repaso del marco teórico}
El problema relativo de dos cuerpos, una vez integradas sus cuadraturas, conduce a la ecuación de la trayectoria en coordenadas polares:
\[
 r(\theta) = \frac{p}{1 + e\cos\theta}
\]
donde \(p = h^2/\mu\) es el parámetro semilatus rectum, \(e = |\vec{e}|\) es la excentricidad y \(\theta\) es la verdadera anomalía medida desde la dirección del vector de Laplace \(\vec{e}\).

Esta expresión describe completamente la forma de la órbita, pero está escrita en un sistema de coordenadas muy particular: aquel cuyo eje polar coincide con la dirección de \(\vec{e}\) y cuyo origen se sitúa en el foco gravitacional. A este marco se le denomina \textbf{sistema natural de la cónica}.

El objetivo analítico de esta clase es formalizar la geometría de las secciones cónicas en su sistema natural, identificar las limitaciones que presenta para la descripción observacional y construir rigurosamente las herramientas de rotación plana necesarias para pasar del sistema natural al sistema de referencia del observador.

\begin{importante}
La ecuación polar de la cónica no depende solo de la forma geométrica de la órbita, sino también de la elección del sistema de referencia. Su forma simple se obtiene únicamente en el sistema natural alineado con la órbita.
\end{importante}

\section{El sistema natural de la cónica}
Toda sección cónica posee una geometría intrínseca que se describe de manera más simple cuando se adopta un sistema de ejes alineado con sus elementos de simetría.

\subsection{Definición del marco natural \texorpdfstring{$(x',y')$}{(x',y')}}
Se define un sistema cartesiano plano \((x', y')\) con las siguientes características:
\begin{itemize}
    \item el eje \(x'\) coincide con el \textbf{eje de simetría} de la cónica,
    \item el origen puede ubicarse en el \textbf{foco} (sistema focal) o en el \textbf{centro geométrico} (sistema central), dependiendo de la conveniencia analítica,
    \item el eje \(y'\) es perpendicular a \(x'\) y completa la base ortonormal derecha del plano orbital.
\end{itemize}

En este sistema, la descripción analítica se simplifica notablemente porque la cónica presenta simetría respecto a \(x'\) y \(y'\).

\subsection{Ecuaciones en el sistema natural}
Dos descripciones típicas en este sistema son:
\begin{itemize}
    \item \textbf{Ecuación polar con origen en el foco:}
\[
 r = \frac{p}{1 + e\cos\theta'}
\]
    donde \(\theta'\) es el ángulo medido desde el semieje positivo \(x'\) hasta el vector posición \(\vec{r}\).
    \item \textbf{Ecuaciones paramétricas cartesianas:} dependiendo del tipo de cónica, las coordenadas \((x', y')\) pueden expresarse en función de un parámetro auxiliar. Por ejemplo, para la elipse usando la anomalía excéntrica \(E\):
\[
 x' = a(\cos E - e), \quad y' = a\sqrt{1-e^2}\sin E
\]
\end{itemize}

Estas expresiones conservan su forma funcional dentro del sistema natural, independientemente de cómo esté orientada la órbita en el espacio físico.

\section{El sistema de referencia del observador y el problema de alineación}
En la práctica astronómica, las posiciones y velocidades se miden respecto a un sistema de referencia fijo o inercial, denominado \textbf{sistema del observador} \((x,y,z)\). Este sistema puede ser, por ejemplo:
\begin{itemize}
    \item el plano ecuatorial terrestre,
    \item el plano de la eclíptica,
    \item el plano tangente a la esfera celeste en la dirección de un objeto.
\end{itemize}

El problema fundamental es que casi nunca el sistema de ejes utilizado en las observaciones coincide con el sistema natural de la órbita. El plano orbital puede estar inclinado respecto al plano de referencia y la línea de ápsides, es decir, la dirección de \(\vec{e}\), puede formar un ángulo arbitrario con el eje \(x\) del observador.

Por lo tanto, para comparar predicciones teóricas con datos observacionales, es necesario establecer una transformación matemática rigurosa que relacione las coordenadas del sistema natural \((x',y')\) con las del sistema del observador \((x,y)\). Como el movimiento sigue siendo plano en esta etapa, la transformación se reduce inicialmente a una \textbf{rotación en el plano}.

\begin{importante}
La dificultad observacional no está en la ecuación de la órbita, sino en que los datos rara vez están expresados en el sistema donde esa ecuación adopta su forma más simple.
\end{importante}

\section{Rotación de la cónica en el plano}
Consideremos dos sistemas cartesianos ortonormales que comparten el mismo origen:
\begin{itemize}
    \item sistema del observador: base \(\{\hat{i}, \hat{j}\}\),
    \item sistema natural de la cónica: base \(\{\hat{i}', \hat{j}'\}\).
\end{itemize}

El sistema natural está rotado un ángulo \(\theta\) respecto al sistema del observador, en sentido antihorario. El vector posición \(\vec{r}\) es el mismo objeto geométrico, pero sus componentes cambian con la base elegida:
\[
\vec{r} = x\hat{i} + y\hat{j} = x'\hat{i}' + y'\hat{j}'
\]

\subsection{Relación entre vectores unitarios}
Para expresar los componentes en una base en función de los de la otra, proyectamos los vectores unitarios rotados sobre la base fija:
\[
\hat{i}' = \cos\theta\,\hat{i} + \sin\theta\,\hat{j}
\]
\[
\hat{j}' = -\sin\theta\,\hat{i} + \cos\theta\,\hat{j}
\]

Sustituyendo en \(\vec{r} = x'\hat{i}' + y'\hat{j}'\) y agrupando términos respecto a \(\hat{i}\) y \(\hat{j}\), obtenemos:
\[
 x\hat{i} + y\hat{j} = (x'\cos\theta - y'\sin\theta)\hat{i} + (x'\sin\theta + y'\cos\theta)\hat{j}
\]

Igualando componentes, se deduce la \textbf{regla de rotación directa}:
\[
\begin{cases}
 x = x'\cos\theta - y'\sin\theta \\
 y = x'\sin\theta + y'\cos\theta
\end{cases}
\]

\section{Derivación analítica de la matriz de rotación \texorpdfstring{$R_z(\theta)$}{Rz(theta)}}
El sistema anterior puede escribirse de manera compacta mediante álgebra matricial. Definimos los vectores columna de coordenadas:
\[
\begin{pmatrix} x \\ y \end{pmatrix} = \mathbf{R}_z(\theta)\begin{pmatrix} x' \\ y' \end{pmatrix}
\]

Aquí, \(\mathbf{R}_z(\theta)\) es la \textbf{matriz de rotación en el plano \(x\)-\(y\) alrededor del eje \(z\)}:
\[
\mathbf{R}_z(\theta) =
\begin{pmatrix}
\cos\theta & -\sin\theta \\
\sin\theta & \cos\theta
\end{pmatrix}
\]

\subsection{Propiedades analíticas de \texorpdfstring{$\mathbf{R}_z(\theta)$}{Rz(theta)}}
Esta matriz posee propiedades fundamentales:
\begin{enumerate}
    \item \textbf{Ortogonalidad:}
\[
\mathbf{R}_z^T(\theta)\mathbf{R}_z(\theta) = \mathbf{I}
\]
    lo cual garantiza que la transformación preserva productos escalares, distancias y ángulos.
    \item \textbf{Determinante unitario:}
\[
\det(\mathbf{R}_z) = \cos^2\theta + \sin^2\theta = 1
\]
    indicando que la rotación es propia y no incluye reflexiones.
    \item \textbf{Preservación de la norma:}
\[
 x^2 + y^2 = x'^2 + y'^2
\]
    consistente con la invariancia de la distancia radial \(r = |\vec{r}|\).
\end{enumerate}

Esta matriz es la herramienta fundamental para trasladar cualquier descripción geométrica de la cónica desde su sistema natural al sistema del observador.

\section{La rotación inversa y el cambio de signo angular}
Hasta ahora se ha trabajado en el caso plano, donde bastaba una sola rotación alrededor del eje \(z\). Sin embargo, para describir una órbita real en un sistema inercial del observador es necesario introducir explícitamente las tres coordenadas \((x,y,z)\) y la orientación espacial completa del plano orbital.

La idea central sigue siendo la misma: el sistema natural de la órbita es el sistema donde la cónica tiene la forma más simple. Lo que cambia es que ahora ese sistema debe incrustarse dentro de un espacio tridimensional.

\subsection{Rotación inversa en el caso plano como antesala del caso espacial}
En dos dimensiones, la transformación inversa se obtenía mediante
\[
\begin{pmatrix} x' \\ y' \end{pmatrix} = \mathbf{R}_z^{-1}(\theta)\begin{pmatrix} x \\ y \end{pmatrix}
\]
y, como \(\mathbf{R}_z\) es ortogonal,
\[
\mathbf{R}_z^{-1}(\theta) = \mathbf{R}_z^T(\theta) = \mathbf{R}_z(-\theta).
\]

Este resultado sigue siendo válido en tres dimensiones para cualquier matriz de rotación elemental: la transformación inversa se obtiene invirtiendo el orden de las rotaciones y cambiando el signo de los ángulos.

\subsection{Unicidad de la transformación en el espacio}
En el problema tridimensional no aparecen matrices distintas para posición y para velocidad. La misma transformación geométrica que orienta el sistema perifocal dentro del espacio inercial actúa sobre cualquier vector del plano orbital. Por ello, una vez construida la matriz total de rotación, esta sirve tanto para transformar \(\vec{r}\) como \(\vec{v}\), o cualquier otro vector expresado en el sistema natural.

\section{Síntesis y preparación para la tridimensionalidad}
Los resultados obtenidos hasta aquí muestran que la descripción plana de la cónica es solamente una etapa intermedia. Para conectar la órbita con el sistema del observador es necesario reconstruir la geometría en tres dimensiones usando coordenadas \((x,y,z)\).

\subsection{7.4 Deducción tridimensional en coordenadas \texorpdfstring{$x,y,z$}{x,y,z}}
Introducimos el sistema natural tridimensional de la órbita, también llamado \textbf{sistema perifocal}, con base ortonormal \(\{\hat{p},\hat{q},\hat{w}\}\):
\begin{itemize}
    \item \(\hat{p}\) apunta hacia el periastro,
    \item \(\hat{q}\) yace en el plano orbital y es perpendicular a \(\hat{p}\),
    \item \(\hat{w}\) es paralelo al momento angular específico \(\vec{h}\).
\end{itemize}

En este sistema, la órbita sigue contenida en un plano, por lo que la coordenada perpendicular vale cero. Así, la posición relativa y la velocidad relativa se escriben como
\[
\vec{r}_{\mathrm{pf}} =
\begin{pmatrix}
 r\cos f \\
 r\sin f \\
 0
\end{pmatrix}
\]
\[
\vec{v}_{\mathrm{pf}} =
\sqrt{\frac{\mu}{p}}
\begin{pmatrix}
 -\sin f \\
 e+\cos f \\
 0
\end{pmatrix}
\]
donde \(f\) es la verdadera anomalía.

El problema consiste en expresar estos vectores en el sistema inercial del observador \(\{\hat{i},\hat{j},\hat{k}\}\), es decir, en términos de las coordenadas \((x,y,z)\). Para ello deben introducirse tres ángulos orbitales:
\begin{enumerate}
    \item la longitud del nodo ascendente \(\Omega\),
    \item la inclinación orbital \(i\),
    \item el argumento del periastro \(\omega\).
\end{enumerate}

La construcción geométrica se hace en tres pasos.

\textbf{Paso 1. Rotación por \(\omega\) alrededor de \(z\).}  
Esta rotación orienta el eje \(\hat{p}\) dentro del plano orbital. La matriz asociada es
\[
\mathbf{R}_z(\omega)=
\begin{pmatrix}
\cos\omega & -\sin\omega & 0 \\
\sin\omega & \cos\omega & 0 \\
0 & 0 & 1
\end{pmatrix}.
\]

\textbf{Paso 2. Rotación por \(i\) alrededor de \(x\).}  
Esta rotación inclina el plano orbital respecto al plano de referencia:
\[
\mathbf{R}_x(i)=
\begin{pmatrix}
1 & 0 & 0 \\
0 & \cos i & -\sin i \\
0 & \sin i & \cos i
\end{pmatrix}.
\]

\textbf{Paso 3. Rotación por \(\Omega\) alrededor de \(z\).}  
Esta última rotación fija la dirección de la línea de nodos en el sistema inercial:
\[
\mathbf{R}_z(\Omega)=
\begin{pmatrix}
\cos\Omega & -\sin\Omega & 0 \\
\sin\Omega & \cos\Omega & 0 \\
0 & 0 & 1
\end{pmatrix}.
\]

La matriz total que transforma del sistema perifocal al sistema inercial es entonces
\[
\mathbf{Q} = \mathbf{R}_z(\Omega)\,\mathbf{R}_x(i)\,\mathbf{R}_z(\omega).
\]
Por tanto,
\[
\vec{r}_{\mathrm{I}} = \mathbf{Q}\,\vec{r}_{\mathrm{pf}},
\qquad
\vec{v}_{\mathrm{I}} = \mathbf{Q}\,\vec{v}_{\mathrm{pf}}.
\]

Al desarrollar el producto matricial y aplicar la matriz a la posición perifocal, se obtiene la deducción explícita en coordenadas cartesianas tridimensionales:
\[
\begin{pmatrix}
 x \\
 y \\
 z
\end{pmatrix}
=
\mathbf{Q}
\begin{pmatrix}
 r\cos f \\
 r\sin f \\
 0
\end{pmatrix}
=
 r
\begin{pmatrix}
 \cos\Omega\cos(\omega+f)-\sin\Omega\sin(\omega+f)\cos i \\
 \sin\Omega\cos(\omega+f)+\cos\Omega\sin(\omega+f)\cos i \\
 \sin(\omega+f)\sin i
\end{pmatrix}.
\]

De aquí se leen directamente las coordenadas espaciales de la órbita:
\[
 x = r\bigl[\cos\Omega\cos(\omega+f)-\sin\Omega\sin(\omega+f)\cos i\bigr]
\]
\[
 y = r\bigl[\sin\Omega\cos(\omega+f)+\cos\Omega\sin(\omega+f)\cos i\bigr]
\]
\[
 z = r\sin(\omega+f)\sin i.
\]

Estas expresiones constituyen la forma explícita de la órbita kepleriana en el espacio tridimensional del observador.

\begin{importante}
La coordenada \(z\) aparece únicamente porque el plano orbital está inclinado. Si \(i=0\), entonces \(z=0\) y se recupera exactamente el caso plano estudiado antes.
\end{importante}

\subsection*{7.5 Velocidad tridimensional en el sistema inercial}
La misma matriz \(\mathbf{Q}\) transforma la velocidad desde el sistema perifocal al sistema inercial. Por tanto,
\[
\begin{pmatrix}
 \dot{x} \\
 \dot{y} \\
 \dot{z}
\end{pmatrix}
=
\mathbf{Q}
\sqrt{\frac{\mu}{p}}
\begin{pmatrix}
 -\sin f \\
 e+\cos f \\
 0
\end{pmatrix}.
\]

Desarrollando componente a componente:
\[
\dot{x} = \sqrt{\frac{\mu}{p}}\left[-\cos\Omega\sin(\omega+f)-\sin\Omega\cos(\omega+f)\cos i + e\bigl(-\cos\Omega\sin\omega-\sin\Omega\cos\omega\cos i\bigr)\right]
\]
\[
\dot{y} = \sqrt{\frac{\mu}{p}}\left[-\sin\Omega\sin(\omega+f)+\cos\Omega\cos(\omega+f)\cos i + e\bigl(-\sin\Omega\sin\omega+\cos\Omega\cos\omega\cos i\bigr)\right]
\]
\[
\dot{z} = \sqrt{\frac{\mu}{p}}\left[\cos(\omega+f)\sin i + e\cos\omega\sin i\right].
\]

Así, la órbita queda completamente descrita en el sistema del observador mediante las seis componentes \((x,y,z,\dot{x},\dot{y},\dot{z})\).

\subsection*{7.6 Interpretación geométrica final}
La dinámica kepleriana sigue siendo plana en su naturaleza: el cuerpo se mueve siempre en un plano fijo perpendicular a \(\vec{h}\). Sin embargo, para un observador inercial, ese plano puede encontrarse arbitrariamente orientado en el espacio. La tridimensionalidad no introduce una nueva fuerza ni una nueva ecuación orbital; introduce una nueva \textbf{geometría de orientación}.

En consecuencia:
\begin{itemize}
    \item la ecuación polar fija la forma de la órbita en su plano natural,
    \item los ángulos \(\Omega\), \(i\) y \(\omega\) fijan la orientación del plano y del periastro en el espacio,
    \item las fórmulas de \((x,y,z)\) y \((\dot{x},\dot{y},\dot{z})\) permiten comparar directamente la teoría con observaciones y simulaciones numéricas.
\end{itemize}

Este es el puente completo entre la solución analítica del problema de dos cuerpos y su representación en un sistema inercial tridimensional.

\begin{importante}
La geometría de las cónicas no solo determina la forma de la trayectoria orbital. También fija la estructura algebraica de las transformaciones de coordenadas que hacen posible comparar teoría, simulación y observación.
\end{importante}